\section{Historical detector studies and component measurements}
\label{app:detector-studies}
\suppressfloats[t]
\subsection{Historical evidence and provenance limits}
We retain the 3,168-trial instruction sweep and related recipient experiments in
Appendix~\ref{app:historical}. Their source traces support arithmetic
reproduction, but some archived backend aliases lack provider-returned version
identifiers. We therefore do not use these rows to establish a model ranking or
a capability--privacy relationship. The fixed-checkpoint structured tasks and
the external task evaluation in Appendix~\ref{sec:external-dojo} provide separately scoped evidence for the
release mechanisms.

\subsection{Comparison with real baselines}
A hand-rolled regular-expression scanner is not a credible DLP baseline, so we run
\textbf{Microsoft Presidio}~\cite{presidio} itself, in two realistic configurations: a
strict identifier policy (SSN, card, email, phone, bank, passport, licence) and a
broad one that adds \texttt{PERSON}, \texttt{LOCATION}, \texttt{NRP} and
\texttt{DATE\_TIME}. The set is 32 realistic scenarios (17 leaks / 15 clean) spanning
all four violation types plus the cases that separate the approaches: non-regex PII, a
diagnosis, authorized external sharing, benign traffic, and single-hint near-misses
that must \emph{not} fire (Table~\ref{tab:frontier}).

The configured baseline observers receive raw content, while our central auditor
receives labels computed locally. This is a deployment choice in this comparison,
not an inherent limitation of Presidio or other tools: they could also run locally.
Our local tagger likewise reads content. The comparison therefore measures these
specific integrations, not superiority over every possible DLP deployment.
\emph{On accuracy}, no Presidio configuration wins on both axes. The identifier policy
reaches only $F1{=}0.26$: it has no notion of owner, purpose, or hop budget, so it
misses inference, purpose, and TTL violations entirely, and the three clean scenarios
it does flag are precisely the \emph{authorized} external shares. The broad policy
buys recall ($0.65$) only by firing indiscriminately---specificity falls to $0.47$,
flagging eight of fifteen clean scenarios including a lunch invitation and a slide
deck---so its apparently strong $83\%$ ``inference recall'' is not inference detection
at all, merely the observation that those messages contain a place and a date. A
policy-aware LLM-judge that reads every byte is precise but insensitive
($F1{=}0.58$, $0\%$ inference recall). On these 32 scenarios, the configured lexical auditor reaches $F1{=}0.94$
with specificity $1.00$ and no raw message text in its central view; the observed
misses concern tagger coverage
(employment- and behavioral-domain hints), never a false alarm. Consistent with the
tagger study below, the open-weight LLM tagger does not help here---it slightly
\emph{trails} the lexical floor---so the lift we previously measured was specific to a
proprietary backend.

\begin{table}[t]\centering\small
\begin{tabular}{lcccccc}
\toprule
Detector & P & R & F1 & spec. & inf.\ rec. & content$\to$center \\
\midrule
Presidio DLP, identifiers & 0.50 & 0.18 & 0.26 & 0.80 & 0\% & 1{,}839 ch \\
Presidio DLP, broad & 0.58 & 0.65 & 0.61 & 0.47 & 83\% & 1{,}839 ch \\
Llama Guard 3 (1B) & 0.53 & 0.53 & 0.53 & 0.47 & 0\% & 1{,}839 ch \\
LLM-judge, policy-aware & 1.00 & 0.41 & 0.58 & 1.00 & 0\% & 1{,}839 ch \\
Ours --- LLM tagger, open 14B & 1.00 & 0.82 & 0.90 & 1.00 & 50\% & \textbf{0} \\
\textbf{Ours --- lexical tagger} & \textbf{1.00} & \textbf{0.88} & \textbf{0.94} & \textbf{1.00} & 67\% & \textbf{0} \\
\bottomrule
\end{tabular}
\caption{Ours vs.\ real baselines on 32 realistic scenarios (17 leaks). All baselines
read every byte; ours decides at zero content egress. Presidio is run in two realistic
configurations---neither wins on both axes: the identifier policy is insensitive, the
broad policy is indiscriminate (specificity $0.47$). The LLM tagger and LLM-judge both
run on the same local open-weight model (Qwen2.5-14B); Llama Guard is the 1B variant
(\S\ref{sec:guard}).}
\label{tab:frontier}
\end{table}

\paragraph{A shipped guardrail.} Presidio is a pattern matcher and the LLM-judge is
a general model handed our policy; neither is a product built to screen agent
messages. Llama Guard is, so we ran it on the message that crosses. It lands at
$F1=0.53$ on this scenario distribution and the two ways it fails are more informative
than the number.\label{sec:guard}

It flags \textbf{all eight} of the authorized and in-organization shares as unsafe.
Those carry exactly the same sensitive content as the leaks; the only difference is
that the recipient is entitled to it. And it calls \textbf{all six} compositional
inference cases safe, because no single message in them contains a hazard.

Both failures are architectural rather than a matter of capacity. The hazard
taxonomy has no slot for \emph{who may receive} a datum, so an authorized share and
a leak are the same message to it; and it judges messages one at a time, so
evidence that only becomes disclosive when it accumulates is invisible. A larger
Llama Guard would share both blind spots exactly --- which is why we report the 1B
variant, forced by local disk, without claiming it is the strongest available: the
limit being demonstrated is not one that scale removes.

\subsection{Adaptive evasion: what can an adversary learn while evading?}
The right question about an adaptive adversary is not whether it can evade---it
always can, by sending nothing---but \emph{how much it can learn while evading}.
We therefore sample the adversary's strategy space at random and run every
strategy through the real auditor, recording both whether the detector fired and
the attribute reported by the recipient under the experiment's scoring rule,
$g=P(A\mid\text{fragments})-p_0$. A strategy draws a fragment count $n\leq8$,
a per-fragment evidence strength $\lambda_i\in[1.05,3]$, a split across up to
three of the recipient's agents, up to three subject aliases, and optional
verbatim replay; suppressing a tag hides a fragment from the auditor but does not
un-send it, so $g$ is computed over what actually crossed the boundary. Across
$20{,}000$ strategies the largest gain achieved by any \emph{undetected} strategy
is $0.29992$ against the policy threshold $\delta{=}0.3$: not one evasion crossed
the line (Table~\ref{tab:evasion}). Within this synthetic model, the detector and simulated adversary share the
same gain function, so complete and correctly grouped metadata aligns their
thresholds. This does not validate the gain model against natural-language
recipients or cover missing observations.

The ablations suppress tags or split subject identities. At tag rate $q=0$,
$53.6\%$ of sampled strategies exceed the simulated gain threshold without an
alert; aliasing produces $16.3\%$ harmful evasions. These are modeled integrity
failures, not measurements of resistance to an actual compromised client.
Our report authentication and canonical identifiers provide implementation
hooks for integrity, but their software implementations do not prove key
confinement, input authenticity, or complete mediation. Rejecting an unauthenticated
report exposes missing evidence; it does not undo an unobserved disclosure.
The enrollment calculation $(1-c)\times53.6\%$ is consequently a conditional
mixture calculation assuming every enrolled sender reports completely and
correctly, not an empirical deployment guarantee.

\begin{table}[t]\centering\small
\begin{tabular}{lccc}
\toprule
Configuration & detect & max $g$ evaded & harmful evasions \\
\midrule
Full design & 0.54 & \textbf{0.29992} & \textbf{0.0\%} \\
$-$ attestation ($q{=}0.4$) & 0.20 & 0.90 & 33.9\% \\
$-$ attestation ($q{=}0$) & 0.00 & 0.90 & 53.6\% \\
$-$ canonical subject ids & 0.38 & 0.90 & 16.3\% \\
\bottomrule
\end{tabular}
\caption{Adaptive evasion, $20{,}000$ random strategies per configuration. A
\emph{harmful} evasion is one that goes undetected \emph{and} gives the recipient
a belief gain $\geq\delta{=}0.3$. Under the full design there are none: the
largest gain any evasion achieves is $0.29992$, just under the threshold. Each
ablated mechanism admits evasion at the full attainable gain ($0.90$).}
\label{tab:evasion}
\end{table}

\paragraph{Attacking the tagger instead.} The study above assumes honest tags;
the tagger is the one component that reads content, so we attack it directly. An
attacking model is given the protected fact and, in the white-box arm, the lexical
tagger's entire 141-term lexicon with the instruction to avoid every term. A
hand-off \emph{evades} only if a reader model from a different family, shown the
hand-off alone, judges the fact conveyed \emph{and} the auditor fires no
content-derived violation---saying nothing is not evasion. Across six attackers (7B open-weight to GPT-4o) and 576 hand-offs, the lexical tagger is evaded $55\%$ of the time black-box and $70\%$ white-box; against the 32B attacker it catches none of 48 white-box hand-offs, against GPT-4o two. The LLM tagger, facing the same hand-offs, is evaded $4\%$ $[1,7]$ black-box and $7\%$ $[2,13]$ white-box (cluster bootstrap over facts), and
black-box evasion does not rise with attacker strength; against the 32B attacker, swapping in a smaller tagger from a third family (8B) raises its evasion from $0\%$ and $6\%$ to $8\%$ and $12\%$, so the figure tracks tagger strength rather than a shared lineage. For the open-weight attackers, avoiding the vocabulary costs the message (conveyance $87\%\!\to\!82\%$); GPT-4o hedges instead, conveying the fact in only $62\%$ of its black-box hand-offs. This
is the thesis on the defender's side---a token list fails against any sender able
to follow ``avoid these words,'' while a tagger that judges meaning does not. Table~\ref{tab:frontier}'s lexical row is therefore the wrong one to deploy against a sender who is trying; the LLM tagger ($F1=0.90$ there) is the configuration the attack leaves standing. Query access is the stronger attack, and we ran it: handed the LLM tagger's verdict and three rewrites, two attackers raise evasion to $16\%$ $[7,26]$---and lower conveyance to $60\%$, because what passes the tagger increasingly fails to say the thing. The tagger is not unbeatable; it is expensive to beat, and the price is the message.

\subsection{What the audit costs, and what actually crosses}\label{app:historical-cost}
Lemma~\ref{lem:blind} says the center's view of a message is a hash plus a
categorical label, so its size cannot depend on the message. That is a claim about
the implementation as much as the mathematics, and it is falsifiable: we sweep
content length over three orders of magnitude and measure the bytes that actually
cross. They are constant at \textbf{418 bytes per message} from 64-byte to
64\,KB messages---$0.0\%$ variation across a $1024\times$ range. This is an empirical fixed-schema size result, not a proof of content secrecy;
metadata values and repeated observations can still disclose information.

Throughput is $32$--$40$k messages/s single-process, and per-message cost grows
only $1.24\times$ from $100$ to $50{,}000$ messages: the detectors are linear in
edges, so an audit scales with traffic rather than with history. We state the
comparison against a content-shipping observer as a crossover rather than a win,
because it is one. At $16$\,KB messages the center receives $39\times$ fewer bytes;
at 64-byte messages it receives $6.5\times$ \emph{more}, since 418 bytes of metadata
then outweighs the message it describes. Lemma~\ref{lem:blind} gives a conditional representation-capacity bound, not a
non-invertibility guarantee. Equality patterns, low-entropy hashes and repeated
metadata may expose information even when the representation is small.

\subsection{The self-authored floor}
On all 48 A2A-MT scenarios the auditor scores $P/R/F1=1.0$ (TP$=18$, FP$=0$,
TN$=30$, FN$=0$) with full per-type coverage and no over-firing on single-hint
clean controls, and holds $1.0$ when the templated text is replaced by
LLM-generated prose---the detectors key on labels, not wording. We report this
first and weight it least: we wrote both the detectors and the scenarios, so it
measures the detectors \emph{given correct labels} and nothing more. Everything
that follows is an attempt to break it. Generating the realistic variant did earn
its keep in one way hand-written tests could not: it exposed TTL fragility under
paraphrasing (fixed by \texttt{provenance\_id}) and an invariant false positive on
prose tokens coinciding with label tokens.

\subsection{The local tagger}
The inference
detector accumulates \texttt{inferred\_categories}; in a real deployment these
must be \emph{produced} from content, not hand-set. A local tagger (lexical,
zero-dependency, optional LLM backend) emits tags to the center. A locally hosted backend keeps raw content on the
agent; a remote API backend also exposes it to that provider and requires an
additional trust assumption. On a labeled set with paraphrased cases, the lexical backend recovers explicit
(sensitive) \texttt{category} at $P/R/F1=1.0$ and \texttt{inferred\_categories}
at $P=1.0$, $R=0.78$; the \textbf{LLM backend (gpt-4o-mini) lifts inferred-category
recall $0.78\to1.0$} (the misses are novel paraphrases like ``the place where they
treat tumors'').
That lift, however, is \emph{backend-dependent}, and we report the negative result:
on the same labeled set, open-weight backends served locally (Qwen2.5-7B/14B,
Llama-3.1-8B) do \emph{not} clear the paraphrase gap---each sits at or below the
lexical floor's $R{=}0.78$, and Llama-3.1-8B is materially worse ($R{=}0.56$)---while
also \emph{costing} explicit-category precision ($1.0\to0.38$--$0.75$) by over-tagging
benign text. Qwen2.5-14B is the partial exception, raising inferred \emph{precision}
to $0.88$ (best inferred $F1{=}0.82$) with recall still at $0.78$. These configuration-dependent results do not isolate capability as a causal
factor; replacing a lexical tagger with a local model does not consistently
improve inference coverage in this experiment.
\texttt{AuditSession.observe} wires it in so a deployer supplies only text plus
policy. Both dishonest suppression and honest labeling errors can remove
necessary evidence. A MAC detects unauthorized report modification, but cannot
distinguish a correctly authenticated labeling error from a correct label.
The following ablations study the semantic recall limits.

\subsection{An author-independent benchmark}
Our 48-scenario benchmark scores 1.0, but we wrote both the detectors and the
scenarios. To reduce direct scenario-author overlap, an LLM generates diverse
scenarios with its own intended labels and we score our auditor against them. Pooling
six batches gives \textbf{75 scenarios} (46 intended leaks), every one of which parses
and is scored---no silent sample shrinkage. The honest result sits far below the
self-authored 1.0: with the lexical tagger $P{=}1.00$, $R{=}0.17$; with an open-weight
LLM tagger recall rises to $R{=}0.46$ but precision falls to $P{=}0.88$ (three false
positives). That last number sharpens a claim we would otherwise have overstated: the
architecture does not over-fire \emph{given correct labels}, but an over-tagging
backend can still manufacture false positives end-to-end, so the absence of false positives in the controlled cases cannot be
generalized to the deployed pipeline or arbitrary correct-label inputs.

\paragraph{Where the recall actually goes.} It would be easy to attribute this gap
to taxonomy coverage and move on; we tested the attribution instead, on a fresh pool
and with the LLM tagger moved to a different model family (Mistral-7B) than the
author, so that one model is not both writing and grading the scenarios---which is
why its absolute numbers differ from the pool above. Scoring one
pool of $99$ author-generated scenarios under progressively more generous taggers
bounds what any tagger fix can buy (Table~\ref{tab:ceiling}), and the result
corrects us in two places.

First, \textbf{the ceiling is $0.76$, not $1.0$}. Even with every part marked
maximally sensitive and carrying every inferred category, roughly a quarter of the
author's intended leaks go undetected. That remainder is not reachable by any
tagger: it is detector scope plus the author's own label noise, and no amount of
taxonomy work will move it. Second, within what \emph{is} reachable, the binding
constraint is the \textbf{sensitivity estimate}, not category coverage---conceding
sensitivity alone recovers the entire reachable gap, while conceding every inferred
category recovers \emph{nothing} ($R{=}0.22$, identical to the lexical floor to the
last confusion-matrix cell). The reason is structural rather than lexical: the
author writes predominantly single-hop scenarios, so converging fragments never
reach $k^{*}{=}2$ and the inference detector cannot contribute at all. Our earlier
reading of the miss taxonomy---that the misses split evenly between disclosure and
inference---was measuring the author's scenario mix, not our detector.

Taxonomy coverage is real but secondary: letting the tagger name the sensitive kind
in open vocabulary instead of choosing from our ten classes lifts recall
$0.37\to0.43$. The oracles are upper bounds and not configurations---flagging
everything costs precision ($0.85\to0.74$)---but they say clearly which component
is worth improving, and that the honest ceiling for this pool is $0.76$ rather than
the $1.0$ a reader would otherwise assume.

\begin{table}[t]\centering\small
\begin{tabular}{lccl}
\toprule
Tagger & $P$ & $R$ & what it isolates \\
\midrule
lexical & 0.85 & 0.22 & the shipped floor \\
LLM backend & 0.86 & 0.37 & shipped, fixed taxonomy \\
LLM, open vocabulary & 0.84 & 0.43 & taxonomy coverage \\
oracle: all inferred categories & 0.85 & 0.22 & category tagging \\
oracle: maximal sensitivity & 0.74 & 0.76 & sensitivity estimate \\
\textbf{oracle: both} & 0.74 & \textbf{0.76} & \textbf{the recall ceiling} \\
\bottomrule
\end{tabular}
\caption{What any tagger fix can buy, on one pool of $99$ author-generated
scenarios. Conceding sensitivity recovers the whole reachable gap; conceding every
inferred category recovers none of it. The ceiling is $0.76$, so a quarter of the
author's intended leaks are out of reach of tagging altogether.}
\label{tab:ceiling}
\end{table}

One caveat keeps this honest in the other direction too: the author model states
some ``inferences'' explicitly and labels borderline cases as leaks, so part of the
$49$--$63\%$ agreement is label noise rather than a detector gap, which is why we
report agreement alongside $P/R$ instead of treating the author's labels as gold.

\subsection{External validation on AgentLeak}
The multi-tenant
auditor reuses a federated engine we validate on the external \textbf{AgentLeak}
benchmark. Driving AgentLeak's own 600 multi-agent scenarios through an LLM to
produce real inter-agent traces and labeling leakage by AgentLeak's rule, the
auditor detects \textbf{100\% of internal-channel leaks at 0.972 precision with
no detected raw vault strings in the central view} ($82.5\%$ internal-channel leak rate,
under the archived scoring rule). On a labeled 33-scenario set the engine holds $P/R/F1=1.0$,
and under the full six-layer desensitizer plus DP it holds $F1\approx0.97$ at
$\varepsilon\in[0.5,3]$ with no detected raw-content matches. These are historical benchmark
measurements, not guarantees of semantic privacy or calibrated inference.

\subsection{Is the operating point cherry-picked?}
The inference detector fires at $k^{*}{=}2$ under our defaults, so it is fair to ask
whether that is a knife-edge fitted to the benchmark. We sweep each knob over the full
48-scenario suite (deterministic; no LLM calls) and find a broad plateau
(Table~\ref{tab:sens}). \emph{Every} setting of $(p_0,\lambda,\delta)$ for which the
closed form gives $k^{*}{=}2$ scores $P/R/F1=1.0$---roughly a $2\times$ band in each
knob independently---and quality degrades only once $k^{*}$ moves off $2$, in the
direction the model predicts: $k^{*}{=}1$ over-fires ($P{=}0.86$) and $k^{*}\ge3$
under-fires ($R{=}0.83$). Within this authored suite, the tested parameter settings track the
closed-form $k^{*}$. This is an internal sensitivity check, not external
calibration of the likelihood ratios or evidence against benchmark-specific tuning. The disclosure floor $\tau$ degrades gracefully rather than cliff-edging
($\tau{=}2$: $P{=}0.90$; $\tau{=}4$: $R{=}0.94$; $\tau{=}5$: $R{=}0.89$), and
raw$\to$center $=0$ at every setting.

\begin{table}[t]\centering\small
\begin{tabular}{lcc}
\toprule
knob & $F1{=}1.0$ band ($k^{*}{=}2$) & degrades at \\
\midrule
$\delta$  & $[0.20,\ 0.40]$ & $0.10$, $0.50$ \\
$\lambda$ & $[2.5,\ 5.0]$   & $2.0$, $9.0$ \\
$p_0$     & $[0.10,\ 0.30]$ & $\le 0.05$ \\
\bottomrule
\end{tabular}
\caption{Inference-detector parameter sensitivity over the 48 scenarios
($\delta$ swept 0.10--0.50, $\lambda$ 1.5--9.0, $p_0$ 0.02--0.30). Quality tracks the
closed-form $k^{*}$, not the raw parameter values.}
\label{tab:sens}
\end{table}


\subsection{Metadata privacy under differential privacy}
Beyond hiding content, we harden the
metadata: identity-bearing fields are pseudonymized with a per-audit shared salt
(cross-tenant comparisons hold in pseudonym space; literal identifiers are
replaced, but linkage and external information may still identify principals) and \texttt{sensitivity} is DP-noised, while categories are kept
structurally. Pseudonymization alone is lossless ($P/R/F1=1.0$). Under DP we
sweep two orders of magnitude, $\varepsilon\in[0.1,8]$, at 40 trials each: mean
$F1$ decays monotonically and gracefully---$0.99$ at $\varepsilon{=}3$, $0.98$ at
$1$, $0.96$ at $0.5$, and still $\approx0.93$ at $0.1$ (sd $\le0.033$)---with no
cliff. The loss is concentrated in recall, i.e.\ disclosure decisions whose noised
sensitivity crosses the floor, exactly where the mechanism predicts, and no raw-content matches are detected at any tested $\varepsilon$.
These accuracy measurements do not establish differential privacy for the whole
report: unchanged categories, routing relations and repeated releases require
separate adjacency and composition analysis.
